\documentclass[a4paper,11pt]{article}

% Sprache/Umlaute
\usepackage[ngerman]{babel}
\usepackage[T1]{fontenc}
\usepackage[utf8]{inputenc}

% Mathematik
\usepackage{amsmath,amssymb}
\usepackage{mathtools}
\usepackage{braket}

%Physik, Chemie
\usepackage{physics}
\usepackage{siunitx}
\sisetup{
  exponent-product = \cdot
}
\usepackage[version=4]{mhchem}
\usepackage{nicefrac}

% Seitenlayout
\usepackage[a4paper]{geometry}
\geometry{
  a4paper,
  left=25mm,
  right=25mm,
  top=25mm,
  bottom=25mm,
}
% -------- Absätze ------------
\setlength{\parskip}{6pt}

\setlength{\emergencystretch}{3em}
\providecommand{\tightlist}{%
  \setlength{\itemsep}{0pt}%
  \setlength{\parskip}{0pt}}

%Stil
\usepackage{enumitem}
\usepackage{titlesec}
\titleformat{\subsection}
{\normalfont}
{}
{0pt}
{}

% Häufige Operatoren
%% Drehimpuls
\newcommand{\Lx}{\hat L_x}
\newcommand{\Ly}{\hat L_y}
\newcommand{\Lz}{\hat L_z}
\newcommand{\Lsq}{\hat L^2}

\newcommand{\px}{\hat p_x}
\newcommand{\py}{\hat p_y}
\newcommand{\pz}{\hat p_z}

\newcommand{\Lp}{\hat L_+}
\newcommand{\Lm}{\hat L_-}

%% Spin
\newcommand{\Sx}{\hat S_x}
\newcommand{\Sy}{\hat S_y}
\newcommand{\Sz}{\hat S_z}
\newcommand{\Ssq}{\hat S^2}

\newcommand{\Sp}{\hat S_+}
\newcommand{\Sm}{\hat S_-}


\newcommand{\ii}{\mathrm{i}}
\newcommand{\hb}{\hbar}

\begin{document}
\section{Allgemeines}

\subsection*{Grundlagen:}
\begin{itemize}[label=--]
    \item $ E = h \nu$
    \item $\lambda \nu = c$
    \item $h = \SI{6.626e-34}{\joule\second}$
    \item $\hbar = \frac{h}{2\pi} =\SI{1.1e-34}{\joule\second}$
    \item $\ce{\unit{\eV} <=>[\cdot \num{1.6e-19}][\colon \num{1.6e-19}] \unit{\joule}}$
\end{itemize}

\subsection*{De Broglie Beziehung:}
\begin{itemize}[label=--]
    \item $p = \frac{h}{\lambda}$
    \begin{itemize}[label= --]
        \item $m=0 \Rightarrow E = pc = h \nu$
        \item $m > 0\Rightarrow E = \underbrace{mc^2 + \frac{p^2}{2m}}_{\text{klassisch}} \underbrace{- \frac{p^4}{8m^3c^2}}_{\text{relativistisch}}$
    \end{itemize}
    \item Der Impuls masseloser Teilchen kann nur indirekt (z.B. über Compont-Experiment) bestimmt werden
    \item Korrespondenzprinzip: Für große Teilchen/Energien muss Übergang von QM in klassische Mechanik erfolgen
\end{itemize}

\subsection*{Schrödingergleichung:}

\begin{itemize}
    \item Zeitunabhängige Schrödingergleichung $ \ii \hb \pdv{}{t}\Psi(x,t) = \hat H \Psi(x,t) $
\end{itemize}

\subsection*{Drehimpuls}

\begin{itemize}
  \item Definition: $\vec{L} = \vec{r} \cross \vec{p}
  = \begin{pmatrix} \hat{x} \\ \hat{y} \\ \hat{z} \end{pmatrix} \cross \begin{pmatrix} \hat{p}_x \\ \hat{p}_y \\ \hat{p}_z \end{pmatrix}
  = \begin{pmatrix} yp_z - zp_y \\ zp_x - xp_z \\ xp_y - yp_x \end{pmatrix}
  = \begin{pmatrix} \hat{L}_x \\ \hat{L}_y \\ \hat{L}_z \end{pmatrix} $ 
  \item Betragsquadrat des Drehimpulses: $\hat{L}^2 = \hat{L}_x^2 + \hat{L}_y^2 + \hat{L}_z^2$
  \item Kommutatoren:
  \begin{itemize}
    \item $[\hat{L}^2, \hat{L}_x] = [\hat{L}^2, \hat{L}_y] = [\hat{L}^2, \hat{L}_z] = 0$
    \item $[\hat{L}_x, \hat{L}_y] = \ii\hb\hat{L}_z \qquad [\hat{L}_z, \hat{L}_x] = \ii\hb\hat{L}_y \qquad [\hat{L}_y, \hat{L}_z] = \ii\hb\hat{L}_x$ \qquad \footnotesize{Immer x,y,z als Reihenfolge}
  \end{itemize}
  \item Gemeinsame Eigenfunktionen von $\hat{L}^2$ und $\hat{L}_z$:
  \begin{itemize}
    \item Kugelflächenfunktionen: $\ket{l, m_l} \equiv Y_{l, m_l}P_{l,m_l}(\cos(\vartheta))\mathrm{e}^{\ii m_l \varphi}$
    \begin{itemize}
      \item $N_{l,m_l}$ ist die Normierungskonstante
      \item $P_{l,m_l}$ ist ein zugeordnetes Legendre-Polynom
    \end{itemize}
    \item Bahndrehimpulsquantenzahl $l$: $l = 0,1,2,3, \dots $ (s-, p-, d-, f- Orbital)
    \item magnetische Bahndreimpuls-QZ $m_l$: $m_l \in {-l, -l+1, \dots, l}$
    \begin{itemize}
      \item $m_l$ gibt die z-Komponente von $l$ an $\rightarrow$ führt bei $l=1$ zu den p-Orbitalen
      \item Entartungsgrad ohne äußeres Magnetfeld: $2l +1$ \footnote{Gibt es ein Symbol für den Entartungsgrad?}
    \end{itemize}
    \item Eigenwerte: $\hat{L}^2\ket{l,m_l} = \hb^2 l(l+1)\ket{l, m_l} \qquad \hat{L}_z \ket{l, m_l} = \hb m_l \ket{l,m_l}$
  \end{itemize}
  \item Auf-/Absteigeoperatoren:
  \begin{itemize}
    \item $\hat{L}_+ = \hat{L}_x + \ii \hat{L}_y$ \quad und \quad $\hat{L}_- = \hat{L}_x - \ii \hat{L}_y$
    \item Wirkung: $\hat{L}_\pm \ket{l,m_l} = \hb \sqrt{l(l+1) - m_l(m_l \color{red}\pm 1 \color{black})}\ket{l, m_l \color{red} \pm 1}$
    \item Verbotene Zustände: $\hat{L}_+ \ket{l, +l} = 0$ \quad und \quad $\hat{L}_-\ket{l,-l}=0$
  \end{itemize}
\end{itemize}

\subsection*{Spin}

\begin{itemize}
\item Eigenschaften analog zum Bahndrehimpuls
\item Betragsquadrat: $\hat S^2 = \hat S_x^2+\hat S_y^2+\hat S_z^2$
\item Kommutatoren:
\begin{itemize}
\item $[\hat S^2,\hat S_x]=[\hat S^2,\hat S_y]=[\hat S^2,\hat S_z]=0$
\item $[\hat S_x,\hat S_y]=\ii\hb\hat S_z \qquad [\hat S_z,\hat S_x]=\ii\hb\hat S_y \qquad [\hat S_y,\hat S_z]=\ii\hb\hat S_x$
\end{itemize}
\item Eigenzustände: $\ket{s,m_s}$
\item Spinquantenzahl: $s=0,\frac12,1,\frac32,\dots$
\begin{itemize}
  \item Fermionen: $s=\frac{1}{2}, \frac{3}{2}, \dots$ $\rightarrow$ halbzahlig
  \item Bosonen: $s = 0,1,2,\dots$ $\rightarrow$ ganzzahlig
\end{itemize}
\item magnetische Spinquantenzahl: $m_s \in {-s,-s+1,\dots,s}$
\item Entartungsgrad: $g=2s+1$
\item Eigenwerte: $\hat S^2\ket{s,m_s}=\hb^2s(s+1)\ket{s,m_s}$ \quad und \quad $\hat S_z\ket{s,m_s}=\hb m_s\ket{s,m_s}$
\item Auf-/Absteigeoperatoren:
\begin{itemize}
\item $\hat S_+=\hat S_x+\ii\hat S_y \qquad \hat S_-=\hat S_x-\ii\hat S_y$
\item $\hat{S}_\pm \ket{s,m_s} = \hb \sqrt{s(s+1) - m_s(m_s \color{red}\pm 1 \color{black})}\ket{s, m_s \color{red} \pm 1}$
\item $\hat S_+\ket{s,s}=0 \qquad \hat S_-\ket{s,-s}=0$
\end{itemize}
\item Elektron: $s=\frac12$, $m_s=\pm\frac12$, $\chi_\pm \equiv \ket{\frac12,\pm\frac12}$
\end{itemize}

\clearpage
\subsection*{Drehimpulskopplung}
\quad $\rightarrow$ gilt für Spin-Bahn oder Spin-Spin Kopplung. Bei Spin-Spin Kopplung gilt: $\hat J \equiv \hat S$ etc. 

\begin{itemize}
\item Gesamtdrehimpuls: $\hat J=\hat J_1+\hat J_2$
\item mögliche Quantenzahlen: $J=j_1+j_2, j_1+j_2-1, \dots,|j_1-j_2|$
\item Projektionen: $m_J=-J,-J+1,\dots,J$
\item Entartungsgrad: $g_J=2J+1$
\begin{itemize}
  \item Kontrolle: $\sum_J(2J+1)=(2j_1+1)(2j_2+1) $
\end{itemize}
\item Spin-Bahn-Kopplung: $j=l+s,\dots,|l-s|$
\item Für Elektronen ($s=\frac12$): $j=l\pm\frac12$
\item Häufigkeiten bzgl. $m_J$: alle möglichen Summen $m_J=m_1+m_2$ bilden; gleiche Summen zählen
\item Operatoren mit Index wirken nur auf den jeweiligen Zustand:

\begin{itemize}
  \item $\hat S_{1}^{2}\ket{s_1,m_1}\ket{s_2,m_2} = \hb^2s_1(s_1+1)\ket{s_1,m_1}\ket{s_2,m_2}$
  \item $\hat S_{+,1}\ket{s_1,m_1}\ket{s_2,m_2} = (\hat S_{+,1}\ket{s_1,m_1})\ket{s_2,m_2} $  
\end{itemize}


\end{itemize}






\end{document}